\documentclass[times, 10pt,twocolumn]{article}
\usepackage{latex8}
\usepackage{times}
\usepackage[utf8]{inputenc}
\usepackage[T1]{fontenc}
\usepackage{booktabs}
\usepackage{url}

\pagestyle{empty}

% Small inline helpers.
\newcommand{\code}[1]{\texttt{\small #1}}
\newcommand{\TODO}[1]{\textbf{[TODO: #1]}}

\begin{document}

\title{dida-trade: A Trading Service Replicated with Vertical Paxos}

\author{Filip Dubj\'ak, n.\textsuperscript{o} 120560 \and
        C\'atia Zharko, n.\textsuperscript{o} 109972 \and
        Joana Cardoso, n.\textsuperscript{o} 109864 \\[4pt]
        Instituto Superior T\'ecnico, Universidade de Lisboa}

\maketitle
\thispagestyle{empty}

\begin{abstract}
\TODO{3--5 sentences once all sections are in: which steps we reached
(1--4), the main design choices (quorum-based phase 1, Multi-Paxos with
pipelined phase 2, causal \code{balance} read from a single replica), and
how each is validated.}
\end{abstract}

%-------------------------------------------------------------------------
\section{Introduction}

dida-trade replicates a small trading state machine over a set of
replicas, using one Paxos instance per log slot. Writes (\code{populate},
\code{add}, \code{sell}, \code{buy}) are totally ordered by consensus and
executed by every replica in slot order. This report covers the steps we
implemented: completing the baseline with a correct phase~1
(Step~1, Section~\ref{sec:step1}), Multi-Paxos (Step~2,
Section~\ref{sec:step2}) \TODO{Step 3 if included}, and a causally
ordered \code{balance} that bypasses Paxos (Step~4,
Section~\ref{sec:step4}). Section~\ref{sec:eval} shows how each property
is checked against the running system.

%-------------------------------------------------------------------------
\section{Step 1: Debug Modes and a Correct Phase 1}
\label{sec:step1}

\subsection{Debug modes}

All debug modes except \code{crash} are implemented by a single gRPC
\code{ServerInterceptor} (\code{DebugInterceptor}) wrapped around the
client-facing and the Paxos services of each replica. The console
service is deliberately \emph{not} intercepted, so a frozen or slow
replica still obeys the console.

\begin{itemize}\itemsep 0pt
\item \textbf{freeze / un-freeze}: an incoming call blocks inside the
  interceptor on a monitor until \code{un-freeze} calls
  \code{notifyAll()}. Requests are therefore delayed, never dropped, and
  resume in arrival order on un-freeze, as the specification requires.
\item \textbf{slow-mode-on/off}: every intercepted call sleeps for a
  random delay in $[500, 4000]$\,ms before being dispatched.
\item \textbf{crash}: the replica first acknowledges the console and then
  calls \code{System.exit} from a separate thread, so the console does not
  see a broken RPC. A crashed replica loses all in-memory state,
  including its Paxos log.
\end{itemize}

Placing the delays in one interceptor, rather than in each RPC handler,
means that every message type (client requests, prepares, accepts,
learns) is affected uniformly. New RPCs added in later steps inherit the
debug behaviour without further code.

\subsection{Why the baseline is unsafe}

The provided phase-1 processor (\code{PhaseOneBogusProcessor}) has three
defects: (i) it reports the ballot as accepted without reading the
acceptor's answer, so a rejection is ignored; (ii) it returns ``done'' on
the \emph{first} reply, so phase~1 never waits for a quorum; and (iii) it
overwrites, rather than compares, the accepted value and ballot of each
reply. Defect (ii) alone breaks safety: a leader that consults one
acceptor that has accepted nothing concludes that no value was chosen,
and proposes a new one for a slot that is already decided.

We reproduce this with the required debug modes (Schedule~A, three
replicas, leader of ballot $b$ is $S_{b \bmod 3}$):

\begin{enumerate}\itemsep 0pt
\item A client issues \code{sell 0 50} (request 101). It is decided in
  slot~0, ballot~0, and executed by $S_0,S_1,S_2$.
\item \code{debug crash 1}; $S_1$ is restarted with an empty log.
\item \code{debug freeze 0}, \code{debug freeze 2}: only the empty $S_1$
  can answer a prepare.
\item \code{ballot 1 1} makes $S_1$ leader; a second client issues
  \code{sell 3 50} (request 102).
\item \code{debug un-freeze} on $S_0$ and $S_2$.
\end{enumerate}

With the baseline, $S_1$ completes phase~1 on its own empty reply,
proposes 102 for slot~0, and executes it. Slot~0 is now \code{101} on
$S_0$ and $S_2$ and \code{102} on $S_1$: the replicas have diverged.

\subsection{Correct phase 1}

\code{PhaseOneResponseProcessor} replaces the baseline. It
(i)~collects promises in a \emph{set} of acceptor ids, so a duplicate
reply cannot fill the quorum twice; (ii)~aborts the ballot on the first
rejection, while remembering the highest ballot seen so the leader can
move past it; and (iii)~for every slot, keeps the accepted value with the
highest \emph{valballot} across the quorum. That value must be
re-proposed in phase~2; only slots with no accepted value are free for
new requests.

One design point deserves mention: \code{getAccepted()} is
\emph{derived} (no rejection \emph{and} $|\textit{promises}| \geq q$)
instead of stored. The response collector never notifies the processor
when an acceptor does not answer, so a stored flag initialised to
\code{true}, as in the baseline, lets a leader proceed on a quorum it
never obtained whenever acceptors are frozen or dead.

In the scenario above, the fixed $S_1$ cannot complete phase~1 while
$S_0$ and $S_2$ are frozen. After un-freeze its quorum includes $S_0$ or
$S_2$, both of which report \code{(valballot 0, value 101)} for slot~0,
so $S_1$ re-proposes 101 and agreement is preserved.

The baseline processor is kept in the code base and selected with
\code{-Ddidatrade.phase1=bogus}. Both processors extend a common
\code{PhaseOneProcessor}, so the main loop is unchanged and the
violation stays reproducible.

\paragraph{Limitation exposed.} After the fix, the restarted $S_1$
agrees that slot~0 holds request 101, but cannot \emph{execute} it: the
consensus value is only a request id, and the request body was sent
before the crash. A restarted replica therefore needs state transfer
(Step~5) to catch up; we do not claim crash recovery.

%-------------------------------------------------------------------------
\section{Step 2: Multi-Paxos}
\label{sec:step2}

\TODO{C\'atia / Joana: multi-prepare over all slots from
\code{next\_log\_entry} while the ballot is unchanged; pipelined phase 2
with in-order execution; recovery of in-flight slots; scheduler changes
for Schedule A/B/C.}

%-------------------------------------------------------------------------
% \section{Step 3: Fast Paxos}
% \TODO{only if implemented}

%-------------------------------------------------------------------------
\section{Step 4: Causally Ordered Balance}
\label{sec:step4}

\subsection{Problem}

In the baseline, \code{balance} is a log entry like any write. Each read
costs a full consensus instance, and reads stall whenever no leader can
make progress (e.g.\ the leader is frozen). Step~4 asks for a
\code{balance} that is not totally ordered, but that respects causal
order and reflects every operation the client has already observed.

\subsection{Design}

Because writes are totally ordered by the Paxos log, a replica's state is
always a \emph{prefix} of the same sequence of commands. Its causal past
is therefore fully described by one integer, the number of log entries
it has executed. We do not need vector clocks.

\begin{itemize}\itemsep 0pt
\item Each replica maintains \code{executed\_index}, the length of the
  executed prefix of the log.
\item Every write reply carries \code{lastindex}, the position of that
  write in the log (plus one). A \code{balance} reply carries the
  replica's \code{executed\_index} when it read the value.
\item The client keeps \code{last\_index}, the maximum \code{lastindex}
  over every reply it has received, and sends it as \code{minindex} in
  each \code{balance} request.
\item A replica receiving \code{balance} waits (on a monitor that the
  main loop signals after each execution) until
  $\code{executed\_index} \geq \code{minindex}$, then reads the
  balance \emph{locally}, without proposing anything.
\end{itemize}

As before, the client sends the request to all replicas and keeps the
first reply. The read therefore completes as soon as \emph{any} replica
that is sufficiently up to date answers.

\subsection{Trade-offs}

\begin{itemize}\itemsep 0pt
\item \textbf{Cost.} A read costs one round trip and no Paxos messages.
  It no longer occupies a log slot and does not depend on a working
  leader: with the leader frozen, \code{balance} is still answered by the
  other replicas.
\item \textbf{Freshness.} Reads may be stale with respect to other
  clients. A linearizable read would need either consensus or leases.
\item \textbf{Lagging replicas.} A replica behind the client blocks the
  request rather than answering with old state. This is safe, and it is
  hidden by the first-reply policy as long as one replica is up to date.
  The wait has no timeout, so a replica that never catches up (e.g.\ the
  restarted replica of Section~\ref{sec:step1}) holds that call open
  indefinitely.
\end{itemize}

%-------------------------------------------------------------------------
\section{Validation and Live Reproduction}
\label{sec:eval}

Every scenario below is a Python script that starts real server
processes, clients and the console, drives them through their terminals
and checks the outcome in their logs. The scripts share one harness
(\code{validation/lab.py}), wait on log events instead of fixed sleeps,
and exit non-zero on any unexpected outcome. Build once with
\code{mvn clean install}.

\subsection{Step 1}

\code{validation/step1/run.py} scripts the scenario of
Section~\ref{sec:step1}:

\begin{small}
\begin{verbatim}
python3 validation/step1/run.py bogus
  -> slot 0: S0=101 S1=102 S2=101  (diverged)
python3 validation/step1/run.py fixed
  -> slot 0: S0=101 S1=101 S2=101  (agreement)
\end{verbatim}
\end{small}

\noindent\textbf{By hand.} Use one terminal per process. Start the
replicas from \code{server/}, the client from \code{app/} and the console
from \code{console/}. To get the bogus run, prefix each server command
with \code{MAVEN\_OPTS="-Ddidatrade.phase1=bogus"}.

\begin{small}
\begin{verbatim}
# server/, i = 0,1,2
mvn exec:java -Dexec.args="9200 i A"
# app/
mvn exec:java -Dexec.args="1 localhost 9200 A"
# console/
mvn exec:java -Dexec.args="localhost 9200 A"

app>     sell 0 50          # slot 0 = 101
console> debug crash 1
         # restart S1 in its terminal
console> debug freeze 0
console> debug freeze 2
console> ballot 1 1
app>     sell 3 50          # request 201
console> debug un-freeze 0
console> debug un-freeze 2
\end{verbatim}
\end{small}

Then compare the line \code{Log entry with number 0 has been decided
with command id = ...} in the three server terminals. With the bogus
processor, $S_1$ prints \code{recovered slots = []} and decides the new
request for slot~0. With the fix, it prints \code{recovered slots = [0]}
and \code{Sending phase 2 for slot 0 in ballot 1, value 101}.

\subsection{Step 4}

\code{python3 validation/step4/run.py} runs the scenario below against
three replicas, a client and the console. To run it by hand, start the processes as in Step~1
(without the bogus flag). Every replica begins with users $0..9$, each
holding 10 currency and 10 items.

\begin{small}
\begin{verbatim}
app>     sell 0 50
app>     buy 1 50
app>     balance 0          # 15
console> debug freeze 0
app>     balance 1          # 5
console> debug un-freeze 0
console> debug freeze 2
app>     sell 2 50
app>     buy 3 50
console> debug un-freeze 2
app>     balance 2          # 15, never 10
\end{verbatim}
\end{small}

\code{balance 0} returns 15, and no server prints \code{Setting response
for command} for it: the read did not go through Paxos. With the leader
$S_0$ frozen, \code{balance 1} is still answered by $S_1$ and $S_2$.
For the last read, $S_2$ has missed both writes. Its terminal prints
\code{minindex 4} and waits until it has executed slots 2 and 3. Then it
prints \code{Balance for reqid 701 is 15 at index 4}. Neither the client
nor $S_2$ ever shows the old balance of 10.

\TODO{Validation of Step 2 (\code{validation/part1}).}

%-------------------------------------------------------------------------
\section{Conclusion}

\TODO{Short summary and known limitations: no state transfer for
restarted replicas; balance wait has no timeout.}

%-------------------------------------------------------------------------
\section*{Authorship Declaration}
\small
We, Filip Dubj\'ak (n.\textsuperscript{o} 120560, MEIC), C\'atia Zharko
(n.\textsuperscript{o} 109972, MEIC) and Joana Cardoso
(n.\textsuperscript{o} 109864, MEIC), hereby declare, on our honor, that
the work entitled ``dida-trade'', carried out as part of the Design and
Implementation of Distributed Applications course, resulted from the
active participation of the following group members. The individual
contributions were:

\begin{itemize}\itemsep 0pt
\item \textbf{Filip Dubj\'ak}: debug modes; demonstration that the
  baseline phase 1 is unsafe and its correction (Step 1); scheduler
  updates for Multi-Paxos; causally ordered \code{balance} (Step 4) and
  its validation.
\item \textbf{C\'atia Zharko}: \TODO{e.g. Multi-Paxos multi-prepare
  across slots.}
\item \textbf{Joana Cardoso}: \TODO{e.g. pipelined phase 2 with ordered
  execution, recovery, review and integration, validation of Steps 1--2.}
\end{itemize}

We declare that this work is original and complies with the applicable
academic regulations.

Lisbon, \TODO{date}

Signatures: \TODO{}

\section*{AI Usage Statement}
All design choices and the code have been by us. The following tools
have been used: Claude and ChatGPT (analysis of the assignment, creation
of tests to verify correct behaviour, support during compilation and
debugging, \TODO{confirm:} drafting of report text); Microsoft Copilot
(clarification of Paxos, Multi-Paxos and Multi-Prepare concepts from the
course materials).

\end{document}
